\documentclass[11pt]{article}
\usepackage[margin=1in]{geometry}
\usepackage{amsmath,amssymb,amsthm,mathtools}
\usepackage{graphicx}
\usepackage{hyperref}

\newtheorem{theorem}{Theorem}
\newtheorem{lemma}[theorem]{Lemma}
\newtheorem{corollary}[theorem]{Corollary}
\newtheorem{remark}[theorem]{Remark}
\newcommand{\Hc}{\mathcal H}
\newcommand{\Rm}{R_{\mathfrak m}}
\newcommand{\C}{\mathbb C}
\newcommand{\N}{\mathbb N}
\newcommand{\linf}{\ell^\infty}

\title{Interpolating zeros force nonseparability in the de Branges gap}
\author{Autonomous solution packet for arXiv:0906.2943}
\date{August 11, 2026}

\begin{document}
\maketitle

\begin{center}
\fbox{\parbox{0.91\textwidth}{
\textbf{Status: candidate partial result, likely valid.}
If the zeros of $\Theta=E^\#/E$ contain a Riesz kernel subsequence of
divergent total imaginary part, then the imaginary-half-line majorant space
contains a bicontinuous copy of $\linf$.  Thus it is neither separable nor
reflexive.  Highly clustered or multiple-zero configurations remain open.
}}
\end{center}

\section{Source question and the unresolved gap}

Let $\Hc=\Hc(E)$ be a de Branges space, put
$\Theta=E^\#/E$, and define on $D=i[1,\infty)$
\[
 \mathfrak m_E(z)=\frac{|E(z)|}{|z+i|}.
\]
The Banach space $\Rm(\Hc)$ consists of all $F\in\Hc$ for which both
$|F|$ and $|F^\#|$ are bounded by a constant multiple of $\mathfrak m_E$
on $D$, with norm
\[
 \|F\|_{\mathfrak m}
 =\max\!\left\{\|F\|_{\Hc},
   \sup_{y\ge1}\frac{|F(iy)|}{\mathfrak m_E(iy)},
   \sup_{y\ge1}\frac{|F^\#(iy)|}{\mathfrak m_E(iy)}\right\}.       \tag{1}
\]

Baranov and Woracek prove that $\Rm(\Hc)=\Hc$ precisely when the mean type
of $\Theta$ is zero and the sum of the imaginary parts of its zeros is
finite.  They also embed $\linf$ when the mean type is negative or the zero
heights have positive limsup \cite[Thms.~3.1--3.2]{BW}.  Remark 3.3 asks
whether the remaining gap can contain a proper $\Rm(\Hc)$ that is separable
or reflexive.

The result below excludes every gap candidate whose zeros carry divergent
height on a Riesz subsequence.  This includes interpolating Blaschke zero
sequences.

\section{Proof intuition}

At a zero $v$ of $\Theta$, the normalized de Branges kernel is, up to a
unimodular constant, $E(z)\sqrt{\operatorname{Im}v}/(\bar v-z)$.  A Cauchy
coefficient $\lambda$ therefore costs $|\lambda|^2/\operatorname{Im}v$ in
squared Hilbert norm.  Equal coefficients are too expensive when the zero
heights tend to zero.

Instead, group a divergent-height Riesz subsequence into remote finite
blocks.  To encode an increment $\Delta_k$, assign to a zero $v$ in block
$k$ the coefficient
\[
 \lambda_v=\Delta_k\frac{\operatorname{Im}v}{H_k},
 \qquad H_k=\sum_{v\text{ in block }k}\operatorname{Im}v.
\]
The coefficients still sum to $\Delta_k$, but their squared Hilbert cost is
exactly $|\Delta_k|^2/H_k$.  Taking $H_k\ge k^4$ makes the total cost finite.
Large gaps between blocks make both relevant Cauchy transforms uniformly
$O(1/y)$; between blocks, their completed sums recover the encoded bounded
sequence.  The formal proof makes these estimates quantitative.

\section{A block Cauchy estimate}

\begin{lemma}\label{lem:block}
Let $V_k$ be finite disjoint sets of points $v=x+ih$ in $\C^+$ with
$h\le1/2$ and $|v|\ge2$.  Write
\[
 P_k=\min_{v\in V_k}|v|,\quad Q_k=\max_{v\in V_k}|v|,
 \quad A_k=\sum_{r\le k}Q_r.
\]
Suppose the blocks are radially ordered and their starts can be chosen so
large that
\[
 P_{k+1}\ge10^6 A_k,                                      \tag{2}
\]
with the same condition imposed recursively after every block is chosen.
For $d=(d_k)\in\linf$, set $d_0=0$, $\Delta_k=d_k-d_{k-1}$, and suppose
coefficients $\lambda_v$ satisfy
\[
 \sum_{v\in V_k}\lambda_v=\Delta_k,qquad
 \sum_{v\in V_k}|\lambda_v|\le2\|d\|_\infty.             \tag{3}
\]
Then
\[
 \sup_{y\ge1}y\left|\sum_{k,v\in V_k}
       \frac{\lambda_v}{\bar v-iy}\right|
 +\sup_{y\ge1}y\left|\sum_{k,v\in V_k}
       \frac{\bar\lambda_v}{v-iy}\right|
 \le C\|d\|_\infty                                      \tag{4}
\]
for an absolute constant $C$.  Moreover, with
$y_k=(A_kP_{k+1})^{1/2}$,
\[
 \left|-iy_k\sum_{r,v\in V_r}
       \frac{\lambda_v}{\bar v-iy_k}-d_k\right|
 \le\frac14\|d\|_\infty.                                \tag{5}
\]
\end{lemma}

\begin{proof}
The elementary denominator bounds are
\[
 |\bar v-iy|\ge\max\{|v|,y\},
 \qquad |v-iy|\ge c\max\{|v|,y\},                        \tag{6}
\]
where the second follows from $h\le1/2$, $y\ge1$, and $|v|\ge2$.
If $P_k\le y<P_{k+1}$, split the sums into completed blocks $r<k$, the
current block, and future blocks.  For completed blocks use
\[
 \frac1{\bar v-iy}=-\frac1{iy}
 +\frac{\bar v}{iy(\bar v-iy)},                            \tag{7}
\]
and the analogous identity with $v$ in place of $\bar v$.
The main term has modulus at most $\|d\|_\infty/y$.  After multiplication
by $y$, the error is bounded by
$2\|d\|_\infty A_{k-1}/y$, hence uniformly by (2).  The entire current
block costs at most a constant times $2\|d\|_\infty/y$ by (3) and (6).
The future blocks cost at most
\[
 2\|d\|_\infty\sum_{r>k}P_r^{-1}\lesssim
 \frac{\|d\|_\infty}{P_{k+1}},
\]
because (2) makes the $P_r$ grow geometrically.  This proves (4) for
$y\ge P_1$.  If $1\le y<P_1$, every block is a future block, so the same
geometric estimate proves (4) directly.

At $y_k$, all blocks through $k$ are completed.  Equation (7) gives the
main term $d_k$.  The past error is at most
$2\|d\|_\infty A_k/y_k$, and the future error after multiplication by
$y_k$ is at most a constant times
$2\|d\|_\infty y_k/P_{k+1}$.  Both are bounded by a fixed constant times
\[
 \|d\|_\infty\sqrt{A_k/P_{k+1}}.
\]
The numerical factor $10^6$ in (2), enlarged during the recursive choice if
needed, makes their sum at most $\|d\|_\infty/4$.  This proves (5).
\end{proof}

\section{Main partial theorem}

\begin{theorem}\label{thm:main}
Let $\Hc(E)$ be a de Branges space and let $(v_n)$ be a subsequence of the
zeros of $\Theta=E^\#/E$ in $\C^+$.  Assume
\begin{enumerate}
\item the normalized kernels
$\widetilde K(v_n,\cdot)=K(v_n,\cdot)/\|K(v_n,\cdot)\|_{\Hc}$
form a Riesz sequence in $\Hc(E)$;
\item $\sum_n\operatorname{Im}v_n=\infty$.
\end{enumerate}
Then $\langle\Rm(\Hc),\|\cdot\|_{\mathfrak m}\rangle$ contains a
bicontinuous copy of $\linf$.  Consequently it is neither separable nor
reflexive.
\end{theorem}

\begin{proof}
Discarding finitely many points, assume
$h_n:=\operatorname{Im}v_n\le1/2$ and $|v_n|\ge2$.  Since every tail still
has divergent height, choose radially ordered finite blocks $V_k$ satisfying
the separation in Lemma \ref{lem:block} and
\[
 H_k:=\sum_{v_n\in V_k}h_n\ge k^4.                          \tag{8}
\]
This is done recursively: skip beyond the required next radius, then take a
finite initial segment of the remaining divergent-height tail.

For $d\in\linf$, define, for $v_n\in V_k$,
\[
 \lambda_n=(d_k-d_{k-1})\frac{h_n}{H_k}.                   \tag{9}
\]
Then (3) holds, and the decisive energy identity is
\[
 \sum_{v_n\in V_k}\frac{|\lambda_n|^2}{h_n}
 =\frac{|d_k-d_{k-1}|^2}{H_k}.                             \tag{10}
\]

At a zero $v_n$ of $\Theta$, the normalized-kernel formula used in the proof
of Proposition 3.8 of \cite{BW} reduces to
\[
 \frac{\widetilde K(v_n,z)}{E(z)}
 =\frac{\sqrt{h_n}}{i\sqrt\pi}\,
   \frac{\overline{E(v_n)}}{|E(v_n)|}\,
   \frac1{\bar v_n-z}.                                    \tag{11}
\]
Choose the coefficient of $\widetilde K(v_n,\cdot)$ so that its quotient by
$E$ in (11) is $\lambda_n/(\bar v_n-z)$, and let $F_d$ be the resulting
series.  The coefficient modulus is
$\sqrt\pi|\lambda_n|/\sqrt{h_n}$.  The Riesz synthesis bound, (8), and
(10) give
\[
 \|F_d\|_{\Hc}^2
 \lesssim\sum_k\frac{|d_k-d_{k-1}|^2}{H_k}
 \le4\|d\|_\infty^2\sum_k k^{-4}.                         \tag{12}
\]
Thus the series converges in $\Hc$ and
\[
 \frac{F_d(z)}{E(z)}=\sum_{k,v_n\in V_k}
       \frac{\lambda_n}{\bar v_n-z}.                      \tag{13}
\]
Taking sharp conjugates gives
\[
 \frac{F_d^\#(z)}{E(z)}
 =\Theta(z)\sum_{k,v_n\in V_k}
       \frac{\bar\lambda_n}{v_n-z}.                       \tag{14}
\]

Since $|\Theta(iy)|\le1$, Lemma \ref{lem:block}, (12)--(14), and compact
point-evaluation bounds on $1\le y\le P_1$ imply
\[
 \|F_d\|_{\mathfrak m}\lesssim\|d\|_\infty.             \tag{15}
\]
Indeed, division by $\mathfrak m_E(iy)=|E(iy)|/(y+1)$ turns the two
majorization terms into $(y+1)$ times the two Cauchy transforms.

Finally, (5) and (13) yield
\[
 \|F_d\|_{\mathfrak m}
 \ge y_k\left|\frac{F_d(iy_k)}{E(iy_k)}\right|
 \ge |d_k|-\frac14\|d\|_\infty.
\]
Taking the supremum over $k$ gives
\[
 \|F_d\|_{\mathfrak m}\ge\frac34\|d\|_\infty.          \tag{16}
\]
Hence $d\mapsto F_d$ is linear, bounded, and bounded below from $\linf$
into $\Rm(\Hc)$.  Its range is closed and bicontinuously isomorphic to
$\linf$, proving the theorem.
\end{proof}

\begin{corollary}\label{cor:interp}
If the zeros of $\Theta$ contain an interpolating Blaschke subsequence of
divergent total imaginary part, then $\Rm(\Hc)$ is neither separable nor
reflexive.  In particular, the zero sequence
\[
 v_n=2^n+\frac{i}{n}\qquad(n\ge1)                          \tag{17}
\]
defines a meromorphic-inner/de Branges example in the source's unresolved
height gap for which $\Rm(\Hc)$ nevertheless contains $\linf$.
\end{corollary}

\begin{proof}
For an interpolating Blaschke sequence, normalized kernels at its zeros form
a Riesz sequence.  In (17), the Blaschke sum
$\sum_n (1/n)/(1+4^n)$ converges, while $\sum_n1/n$ diverges and the heights
tend to zero.  Uniform pseudohyperbolic separation follows from the
exponential spacing of the real parts: for $m\ne n$,
\[
 1-\left|\frac{v_n-v_m}{v_n-\bar v_m}\right|^2
 =\frac{4\,\operatorname{Im}v_n\operatorname{Im}v_m}
 {|v_n-\bar v_m|^2},
\]
and the sum of these defects tends to zero as $n\to\infty$; the finitely many
initial products are positive.  Thus the sequence is interpolating.  Apply
Theorem \ref{thm:main}.  The standard correspondence between meromorphic
inner functions and Hermite--Biehler functions supplies $E$ with
$E^\#/E=\Theta$.
\end{proof}

\section{Scope, verification, and novelty}

This theorem closes a substantial part of Remark 3.3's gap but not all of
it.  A divergent zero-height sum can be concentrated in increasingly tight
clusters or high multiplicities.  Such a configuration need not visibly
contain a Riesz subsequence retaining divergent height.  Extending the block
argument would require uniform frame estimates for derivative-kernel blocks;
none is proved here.

The script \texttt{code/verify\_weighted\_blocks.py} checks the exact energy
identity (10), both finite Cauchy-transform bounds, and recovery between five
lacunary blocks over 80 deterministic random trials.  It is a sanity check,
not part of the proof.

A bounded novelty search on August 11, 2026 used the exact source question,
the paper title, ``interpolating zeros,'' ``Riesz kernels,'' and
``nonseparable majorant space.''  It found the Baranov--Woracek series and
adjacent de Branges interpolation literature, but no matching theorem or full
resolution.  Novelty confidence is moderate, not exhaustive.

\begin{thebibliography}{9}
\bibitem{BW}
A. Baranov and H. Woracek,
\emph{Majorization in de Branges spaces II. Banach spaces generated by
majorants}, Collect. Math. 62 (2011), 27--55; arXiv:0906.2943.
\end{thebibliography}

\section*{Source evidence}

Page 7 of arXiv:0906.2943 states the equality criterion, the known
$\linf$-embedding regime, and the separability/reflexivity question in
Remark 3.3.

\begin{center}
\includegraphics[width=0.93\textwidth]{figures/open_problem_crop.png}
\end{center}

\end{document}
